%! Sinusoidal Functions — Unit 6, Lesson 6.6 — Lesson Plan
\documentclass[10pt]{article}
\usepackage{precalculus-article}
\usepackage{precalculus-boxes}
\usepackage{graphicx}
\graphicspath{{images/}}
\newcommand{\TallMath}[1]{$\displaystyle #1\rule[-1.4em]{0pt}{3.2em}$}
\newcommand{\CourseName}{Precalculus}
\newcommand{\MeetingLength}{60 minutes}
\newcommand{\UnitNumberName}{Unit 6: Trigonometric Functions \quad}
\newcommand{\LessonNumberName}{Lesson 6.6: Sinusoidal Functions}

\begin{document}

%-----------------------------------------------------------------------
% 1. TITLE BLOCK
%-----------------------------------------------------------------------
\begin{center}
  {\Large\bfseries\color{navy} Precalculus}\\[4pt]
  {\large\bfseries Unit 6: Trigonometric Functions}\\[2pt]
  {\large\bfseries Lesson 6.6: Sinusoidal Functions}\\[2pt]
  {\normalsize Instructional Periods: 2--3 \quad|\quad Meeting Length: 60 minutes}
\end{center}

\vspace{4pt}
\hrule
\vspace{8pt}

%-----------------------------------------------------------------------
% 2. PRIMARY OBJECTIVE
%-----------------------------------------------------------------------
\begin{tcolorbox}[colback=sky, colframe=navy, boxrule=0.9pt, arc=2mm,
    left=2mm, right=2mm, top=1.5mm, bottom=1.5mm]
  \textbf{Primary Objective} \hfill \textit{Big Idea: Functions}

  \medskip
  Students identify and describe the key characteristics --- period, frequency,
  amplitude, and midline --- of sinusoidal functions, including the symmetry
  properties of sine (odd) and cosine (even).
\end{tcolorbox}

\vspace{6pt}

%-----------------------------------------------------------------------
% 3. PRIORITY IDEAS & SKILLS
%-----------------------------------------------------------------------
\begin{skillbox}[Priority Ideas \& Skills]{goldbox}
  \textbf{Skills:} 2.A (Identify key characteristics from representations) \quad
  3.A (Describe characteristics of functions)

  \medskip
  \textbf{Essential Knowledge:}
  \begin{itemize}
    \item \textbf{EK 3.5.A.1} A \textit{sinusoidal function} involves additive and
          multiplicative transformations of $f(\theta)=\sin\theta$.
          Both sine and cosine are sinusoidal; $\cos\theta = \sin(\theta+\pi/2)$.
    \item \textbf{EK 3.5.A.2} Period and frequency are reciprocals.
          The period of $\sin\theta$ and $\cos\theta$ is $2\pi$; frequency is
          $\tfrac{1}{2\pi}$.
    \item \textbf{EK 3.5.A.3} Amplitude $=\tfrac{\text{max}-\text{min}}{2}$.
          For $\sin\theta$ and $\cos\theta$, amplitude $= 1$.
    \item \textbf{EK 3.5.A.4} Midline $= \tfrac{\text{max}+\text{min}}{2}$.
          For $\sin\theta$ and $\cos\theta$, midline is $y=0$.
    \item \textbf{EK 3.5.A.5} As inputs increase, sinusoidal graphs oscillate
          between concave down and concave up.
    \item \textbf{EK 3.5.A.6} $y=\sin\theta$ has rotational symmetry about the
          origin (odd); $y=\cos\theta$ has reflective symmetry over the $y$-axis
          (even).
  \end{itemize}
\end{skillbox}

%-----------------------------------------------------------------------
% 4. VOCABULARY
%-----------------------------------------------------------------------
\begin{skillbox}[{Vocabulary, Concepts, \& Theorems}]{greenbox}
  \begin{multicols}{3}
  \begin{itemize}
    \item sinusoidal function
    \item period
    \item frequency
    \item amplitude
    \item midline
    \item odd function
    \item even function
    \item concave up
    \item concave down
  \end{itemize}
  \end{multicols}
\end{skillbox}

%-----------------------------------------------------------------------
% 5. ACTIVATE PRIOR KNOWLEDGE
%-----------------------------------------------------------------------
\begin{skillbox}[Activate Prior Knowledge \& Spiral Review]{sky}
  \begin{itemize}
    \item \textbf{Lesson 3.4} --- Graphs of $\sin\theta$ and $\cos\theta$: shape,
          key points (zeros, maxima, minima), and the unit circle connection.
    \item \textbf{Lesson 3.1} --- Periodic functions and period: the smallest
          positive input shift $k$ such that $f(\theta+k)=f(\theta)$ for all $\theta$.
    \item \textbf{Unit 1} --- Even and odd functions: $f(-x)=f(x)$ (even) vs.\
          $f(-x)=-f(x)$ (odd); graphical interpretations (y-axis symmetry vs.\
          rotational symmetry about the origin).
  \end{itemize}
\end{skillbox}

%-----------------------------------------------------------------------
% 6. HOOK
%-----------------------------------------------------------------------
\begin{skillbox}[Hook]{sky}
  \textit{``A recording engineer notices that a sound wave can be described by a
  function with amplitude 0.3 and frequency 440 Hz. What information does that
  tell us about the wave, and how does it relate to the sine function?''}

  \medskip
  \begin{itemize}
    \item Amplitude 0.3 means the wave oscillates 0.3 units above and below its
          centerline --- controlling perceived loudness.
    \item Frequency 440 Hz means 440 complete cycles per second (the note A above
          middle C).
    \item The period of one cycle is $T = \tfrac{1}{440} \approx 0.00227$ seconds.
    \item \textbf{Key idea:} period and frequency are \textit{reciprocals}:
          $\text{frequency} = \dfrac{1}{\text{period}}$.
  \end{itemize}

  \smallskip
  \textit{Transition:} Today we identify the formal characteristics --- amplitude,
  midline, period, and frequency --- that describe any sinusoidal function.
\end{skillbox}

%-----------------------------------------------------------------------
% 7. LESSON
%-----------------------------------------------------------------------
\begin{skillbox}[Lesson]{sky}
  \begin{multicols}{2}
  \textbf{Sinusoidal Functions}

  A \textbf{sinusoidal function} is any function that can be written as an
  additive and/or multiplicative transformation of $f(\theta)=\sin\theta$.
  Both sine and cosine are sinusoidal:
  \[
    \cos\theta = \sin\!\left(\theta + \tfrac{\pi}{2}\right)
  \]

  \columnbreak

  \textbf{Amplitude}

  Half the total vertical range:
  \[
    \text{amplitude} = \frac{\text{max} - \text{min}}{2}
  \]
  Always non-negative.  For $\sin\theta$ and $\cos\theta$: amplitude $=1$.
  \end{multicols}

  \begin{multicols}{2}
  \textbf{Midline}

  Horizontal line halfway between max and min:
  \[
    \text{midline: } y = \frac{\text{max} + \text{min}}{2}
  \]
  For $\sin\theta$ and $\cos\theta$: midline is $y=0$.

  \columnbreak

  \textbf{Period \& Frequency}

  Period $T$: smallest positive value with $f(\theta+T)=f(\theta)$.
  \[
    \text{frequency} = \frac{1}{T}
  \]
  For $\sin\theta$ and $\cos\theta$: $T=2\pi$, frequency $=\tfrac{1}{2\pi}$.
  \end{multicols}

  \textbf{Symmetry}
  \begin{itemize}
    \item $y=\sin\theta$ is \textbf{odd}: $\sin(-\theta)=-\sin\theta$.
          Graph has rotational symmetry about the origin.
    \item $y=\cos\theta$ is \textbf{even}: $\cos(-\theta)=\cos\theta$.
          Graph has reflective symmetry over the $y$-axis.
  \end{itemize}

  \textbf{Concavity:} As $\theta$ increases, sinusoidal graphs oscillate between
  concave down (near maxima) and concave up (near minima).
\end{skillbox}

%-----------------------------------------------------------------------
% 8. EXPLICIT INSTRUCTION
%-----------------------------------------------------------------------
\begin{skillbox}[Explicit Instruction: Identifying Characteristics from a Description]{sky}
  \textbf{Steps:}
  \begin{enumerate}
    \item \textbf{Find max and min.}  Compute:
          amplitude $= \dfrac{\text{max}-\text{min}}{2}$ \quad and \quad
          midline $= \dfrac{\text{max}+\text{min}}{2}$ (write as $y=\ldots$).
    \item \textbf{Determine the period.}
          Read from the graph (length of one complete cycle) or context.
    \item \textbf{Compute frequency.}  Frequency $= \dfrac{1}{\text{period}}$.
  \end{enumerate}

  \textbf{Worked Example:} A sinusoidal function has maximum value $8$ and minimum
  value $-2$, with period $4\pi$.
  \begin{enumerate}
    \item Amplitude $= \dfrac{8-(-2)}{2} = 5$
    \item Midline: $y = \dfrac{8+(-2)}{2} = 3$
    \item Period $= 4\pi$ (given)
    \item Frequency $= \dfrac{1}{4\pi}$
  \end{enumerate}
\end{skillbox}

%-----------------------------------------------------------------------
% 9. ACTIVE MONITORING
%-----------------------------------------------------------------------
\begin{skillbox}[Active Monitoring]{redbox}
  Circulate as students work on guided notes and the group activity.  Watch for:
  \begin{itemize}
    \item Confusing amplitude with max value (amplitude is \textit{half} the range).
    \item Computing midline as average of amplitude and zero rather than of max and min.
    \item Inverting the period/frequency relationship.
    \item Misclassifying sine as even or cosine as odd.
  \end{itemize}

  \textbf{Cold-call check:} ``A function has max $= 5$ and min $= -1$.  What is
  the amplitude?  The midline?'' \quad (amplitude $= 3$; midline $y=2$.)
\end{skillbox}

%-----------------------------------------------------------------------
% 10. GROUP WORK & DIFFERENTIATION
%-----------------------------------------------------------------------
\begin{skillbox}[Group Work \& Differentiation]{redbox}
  Students work in assigned tiers on the Group Activity (separate document).
  \begin{itemize}
    \item \textbf{Tier R (Remediate):} Identifying amplitude, midline, period, and
          frequency from verbal descriptions with max and min given.
          Scaffolded formula prompts are included.
    \item \textbf{Tier A (At Level):} Determining all four characteristics from
          tables of values; classifying symmetry (odd/even/neither) from a table.
    \item \textbf{Tier E (Extend):} Reversing the process (given characteristics,
          find max/min); algebraic justification using the unit circle definition
          that $\sin(\pi-\theta)=\sin\theta$.
  \end{itemize}
  \textbf{Grouping note:} Pairs work best; groups of 3 are acceptable.
\end{skillbox}

%-----------------------------------------------------------------------
% 11. INDIVIDUAL WORK & ASSESSMENT
%-----------------------------------------------------------------------
\begin{skillbox}[Individual Work \& Assessment]{redbox}
  \begin{itemize}
    \item \textbf{Exit Ticket} (last 5 min): 2 items assessing amplitude/midline
          computation and period-to-frequency conversion.
    \item \textbf{Homework}: 5--6 problems including an AP-style multiple-choice
          item and an extension on odd/even symmetry.
    \item \textbf{Formative data:} Collect exit tickets to identify students who
          need reteaching before Lesson 3.6.
  \end{itemize}
\end{skillbox}

%-----------------------------------------------------------------------
% 12. REINFORCEMENT & EXTENSION
%-----------------------------------------------------------------------
\begin{skillbox}[Reinforcement \& Extension]{goldbox}
  \textbf{Reinforce:} Students who struggle with the formulas benefit from returning
  to the unit circle: the maximum $y$-value of $\sin\theta$ is $1$ (at
  $\theta=\pi/2$) and the minimum is $-1$ (at $\theta=3\pi/2$).
  Amplitude $= \tfrac{1-(-1)}{2} = 1$.

  \medskip
  \textbf{Preview --- Lesson 3.6: Sinusoidal Function Transformations}

  In the next lesson, students analyze how parameters $a$, $b$, $c$, and $d$
  transform $f(\theta)=\sin\theta$ in the general form
  $g(\theta) = a\sin(b(\theta+c))+d$, connecting those parameters directly to
  the amplitude ($|a|$), period ($\tfrac{2\pi}{|b|}$), horizontal shift ($-c$),
  and midline ($d$) studied today.
\end{skillbox}

\end{document}
